
Alignment methods for code generation models typically optimize for single metrics such as pass@k correctness, yet may create systematic biases across other quality dimensions including complexity and efficiency. This work develops a diagnostic framework for detecting alignment method signatures---patterns in multi-dimensional performance space that reveal implicit optimization objectives. The methodology measures models across correctness (test execution), complexity (cyclomatic complexity, AST depth), and efficiency (runtime, memory) dimensions, applying PCA-based clustering with Cohen's d effect size analysis. A proof-of-concept validation using simulated performance data demonstrates that the methodology detects simulated signatures when present (simulated Cohen's d=7.835). A minimal real-model validation (3 models, 10 tasks, 300 code generations) confirms the framework's capability to profile actual models, yielding correctness rates of 13.0\% (phi-2), 1.0\% (codegen-350M-mono), and 0.0\% (codegen-350M-nl). The real-model data show the execution-aligned model (phi-2, 2.7B parameters) achieving higher correctness than smaller models, though model scale confounds prevent isolating alignment effects from capacity effects. This proof-of-concept establishes methodological feasibility for signature detection, with real-model validation at scale (164 tasks, matched-scale comparisons) required before claims about alignment method behavior can be substantiated.
